\begin{abstract}
Sending one unpadded integer to multiple RSA recipients with public exponent $e=3$ can expose that integer without any private key. BroadcastGuard studies the exact CRT-and-cube-root condition, compares the same attack against RSA-OAEP under paired keys, and implements authenticated file sharing for several recipients. Twenty independent RSA key-set blocks produced 560 controlled subset observations, all consistent with the predicted integer-recovery condition. In 60 same-key comparisons, the tested attack recovered the textbook RSA message in 60 cases and the OAEP message in none; intended OAEP recipients decrypted in all 60. For a 1\,MiB hybrid-encrypted file and one recipient, JSON/Base64 serialization accounted for 93.78\% of the measured encryption workflow mean on one Windows machine. Eleven selected package-validity cases matched expected outcomes. The contribution is a reproducible educational integration and evidence-based analysis of established cryptographic methods, not a new primitive or a proof of universal security.
\end{abstract}

\section{Introduction and research questions}
RSA public encryption transforms an integer under a public modulus and exponent~\cite{rivest1978method}. If the same integer is encrypted without padding under distinct moduli with a small exponent, the ciphertexts retain an algebraic relationship accessible to an observer. Håstad examined broader systems of low-degree modular equations~\cite{hastad1988solving}; the fixed-message case studied here has a direct Chinese Remainder Theorem (CRT) derivation. This deliberately unsafe example exposes why correct arithmetic and correct security assumptions must be considered together. It does not imply that a properly configured RSA-OAEP deployment is vulnerable to the same simple procedure.

A demonstration that Eve sometimes recovers a message is incomplete without specifying \emph{when} recovery is mathematically possible, what Eve actually learns, and how the outcome changes when the encryption scheme is replaced. It also leaves unanswered how to protect a file too large for direct RSA-OAEP encryption. BroadcastGuard connects these questions through a public-only attack engine, a seven-subset conditions analyzer, a same-key OAEP comparison, and an AES-GCM/RSA-OAEP multi-recipient package. The \href{https://broadcastguard-crypto.streamlit.app/}{interactive application} illustrates the algorithms; its interface is not independent experimental evidence.

The study addresses four questions: \textbf{RQ1}, which numerical condition characterizes recovery of a shared integer by CRT-and-root? \textbf{RQ2}, what changes when only textbook RSA is replaced by randomized RSA-OAEP while keys, plaintext and attacker procedure remain fixed? \textbf{RQ3}, how do file size, recipient count and JSON/Base64 packaging affect hybrid-encryption time and storage? \textbf{RQ4}, which selected invalid or altered packages are rejected without returning plaintext? The group's contribution is the controlled system design, implementation, measurement and interpretation. RSA, CRT, OAEP and GCM are established techniques.
